\documentclass[11pt]{article}
\usepackage[a4paper,margin=1in]{geometry}
\usepackage{amsmath,amssymb,amsthm,mathtools}
\usepackage{booktabs,longtable,array,tabularx}
\usepackage{enumitem}
\usepackage{xcolor}
\usepackage{hyperref}
\usepackage{microtype}
\usepackage{titlesec}
\usepackage{fancyhdr}
\usepackage{lastpage}
\usepackage{listings}

\definecolor{accent}{HTML}{1F4E79}
\definecolor{soft}{HTML}{EEF5FB}
\definecolor{dark}{HTML}{1F2937}
\hypersetup{colorlinks=true,linkcolor=accent,citecolor=accent,urlcolor=accent}
\titleformat{\section}{\Large\bfseries\color{accent}}{\thesection}{0.7em}{}
\titleformat{\subsection}{\large\bfseries\color{dark}}{\thesubsection}{0.7em}{}
\pagestyle{fancy}
\fancyhf{}
\lhead{Referral Auctions on Networks}
\rhead{Abhi Jain, Harshit Raj, Ridam Jan}
\cfoot{\thepage/\pageref{LastPage}}

\newtheorem{definition}{Definition}
\newtheorem{proposition}{Proposition}
\newtheorem{researchq}{Research Question}
\newtheorem{conjecture}{Conjecture}

\newcommand{\N}{\mathcal{N}}
\newcommand{\A}{\mathcal{A}}
\newcommand{\R}{\mathbb{R}}
\newcommand{\eps}{\varepsilon}
\newcommand{\argmax}{\operatorname*{arg\,max}}
\newcommand{\Rev}{\operatorname{Rev}}
\newcommand{\SW}{\operatorname{SW}}

\begin{document}

\begin{titlepage}
    \centering
    \vspace*{1.3cm}
    {\Huge \bfseries Incentive-Compatible Referral Auctions on Networks\\[0.2cm]}
    {\Large A Theory-First Project Report and Experiment Plan\\[0.8cm]}
    \vspace{0.5cm}
    \begin{tabular}{rl}
        Team: & Abhi Jain (23b0903), Harshit Raj (23b0950), Ridam Jan (23b1049)\\
        Theme: & Mechanism design in social networks and referral auctions\\
        Main base paper: & Li, Hao, Zhao, Zhou, \emph{Mechanism Design in Social Networks}\\
    \end{tabular}
    \vfill
    \fbox{\begin{minipage}{0.86\textwidth}
    \textbf{Core research stance.} We do not treat the project as simply implementing known mechanisms. The implementation is a controlled environment for testing theoretical hypotheses around robustness, non-i.i.d. values, topology, and reward/revenue tradeoffs. The strongest candidate contribution is a small theorem or counterexample in one restricted but nontrivial setting, supported by systematic simulations.
    \end{minipage}}
    \vfill
    {\large \today}
\end{titlepage}

\tableofcontents
\newpage

\section{Executive Summary}

Classical auction theory assumes that the seller can directly reach all bidders. Diffusion auctions change the model: the seller initially reaches only her neighbours in a social network, and informed agents strategically decide whether to invite their neighbours. This creates a conflict: the seller wants diffusion, while a bidder may lose the item if she invites a stronger competitor.

The first mechanism in the area, the Information Diffusion Mechanism (IDM) of Li et al.~\cite{li2017mdsn}, uses diffusion-critical nodes to reward useful propagation. Its central message is simple but deep: efficiency, incentive compatibility, and seller revenue cannot all be obtained by a naive extension of VCG. Network-VCG is diffusion incentive-compatible but may run a deficit; IDM restores nonnegative revenue by sacrificing some allocation efficiency.

Recent work extends this in many directions: multi-item sales~\cite{zhao2018multi}, redistribution~\cite{zhang2020redistribution,gu2023redistribution}, fair rewards~\cite{zhang2020fair}, transaction costs~\cite{li2024cost}, privacy~\cite{jia2023dp}, Sybil attacks~\cite{chen2023sybil}, and Myerson-style optimal referral auctions~\cite{bhattacharyya2024optimal}. The most attractive open-ish direction for this course project is not to reprove these papers, but to study a narrow, credible gap:

\begin{quote}
\textbf{How robust are revenue-optimal or high-revenue referral auctions when the i.i.d. valuation assumption is relaxed, especially on trees or nearly-tree-like networks?}
\end{quote}

This is promising because Bhattacharyya et al.~explicitly obtain revenue-optimal structure for i.i.d. bidders and state that the non-i.i.d. setting is far less structured. It is also implementable: we can compare IDM, network-VCG, local Vickrey, and tunable referral mechanisms under controlled topology and valuation models.

\section{Model}

Let $G=(V,E)$ be a graph with seller $s\in V$ and buyers $N=V\setminus\{s\}$. Buyer $i$ has private value $v_i\ge 0$ for a single item. The seller initially informs only $r_s$, her neighbours. A buyer who receives the information submits two reports:
\[
    a_i=(b_i,r_i'), \qquad r_i'\subseteq r_i,
\]
where $b_i$ is her bid and $r_i'$ is the set of neighbours she invites.

An action profile induces a directed invitation graph. A buyer participates only if she is reachable from $s$ in this invitation graph. A mechanism $M=(\pi,p)$ chooses an allocation $\pi_i(a)\in\{0,1\}$ and a payment $p_i(a)\in\R$. Positive $p_i$ means buyer $i$ pays the seller; negative $p_i$ means buyer $i$ is rewarded.

\begin{definition}[Individual rationality]
A mechanism is individually rational (IR) if every buyer obtains nonnegative utility when she reports her value truthfully, irrespective of the actions of others.
\end{definition}

\begin{definition}[Diffusion incentive compatibility]
A mechanism is diffusion incentive-compatible (DIC) if for every buyer, reporting the true value and inviting all neighbours is a dominant strategy.
\end{definition}

\begin{definition}[Weak budget balance]
A mechanism is weakly budget balanced (WBB) if
\[
    \Rev(M,a)=\sum_{i\in N}p_i(a)\ge 0
\]
for every feasible action profile $a$.
\end{definition}

\subsection{Diffusion-critical sets}

Given a feasible invitation graph and a participant $j$, buyer $i$ is a diffusion-critical node of $j$ if every directed path from $s$ to $j$ passes through $i$. Define
\[
    d_i=\{i\}\cup \{j: i\text{ is diffusion-critical for }j\}.
\]
The diffusion-critical sequence $C_m$ of a buyer $m$ is the ordered chain of critical nodes of $m$, followed by $m$ itself. This sequence is the central combinatorial object in IDM.

\section{Known Mechanisms and What They Teach}

\subsection{Central Vickrey auction}

If all bidders are known, the seller allocates to the highest bidder and charges the second-highest bid. This is truthful and efficient, but it is not a network mechanism: it ignores the information propagation problem.

\subsection{Local Vickrey auction without diffusion incentives}

The seller runs a second-price auction only among $r_s$. This is safe but local. It can miss high-value buyers deeper in the network, so welfare and revenue can be much smaller than what would be possible after diffusion.

\subsection{Network VCG}

Network VCG allocates to the highest reachable bidder. The payment of buyer $i$ is the externality she imposes:
\[
    p_i^{VCG}=W(a_{-d_i})-\left(W(a)-\pi_i(a)b_i\right).
\]
For a single item, $W(a)$ is the maximum bid among participants. The key issue is that a critical ancestor of the winner may receive a reward equal to the welfare she enables. On a line graph where only the last node has value $1$, every intermediate node must be rewarded, so the seller revenue becomes negative.

\begin{proposition}[Known]
Network VCG is IR and diffusion incentive-compatible, but it is not weakly budget balanced.
\end{proposition}

\subsection{Information Diffusion Mechanism}

Let $m$ be the highest bidder and $C_m=(c_1,c_2,\ldots,m)$ be its diffusion-critical sequence. IDM allocates to the first node $c_t$ on this sequence whose bid is high enough to win after the next critical descendant block is removed; otherwise $m$ wins. If $w$ wins, then for $i\in C_w\setminus\{w\}$,
\[
    p_i^{IDM}=v^*_{-d_i}-v^*_{-d_{i+1}},
\]
and the winner pays
\[
    p_w^{IDM}=v^*_{-d_w}.
\]
Since $d_i\supseteq d_{i+1}$, the on-path payments are nonpositive, i.e., rewards.

\begin{proposition}[Known]
IDM is IR, diffusion incentive-compatible, and weakly budget balanced. Its revenue is always at least the revenue of network VCG.
\end{proposition}

\subsection{Critical Diffusion Mechanisms and later variants}

Later work generalizes IDM into broader classes of diffusion mechanisms. The high-level pattern is: give earlier critical nodes priority, compensate them when their diffusion causes them to lose, and choose payments through critical thresholds. This opens a design space where one can trade efficiency, revenue, reward fairness, privacy, and robustness.

\section{Literature Map}

\begin{longtable}{p{0.19\textwidth}p{0.25\textwidth}p{0.24\textwidth}p{0.22\textwidth}}
\toprule
\textbf{Paper} & \textbf{Main contribution} & \textbf{Assumptions} & \textbf{Project-relevant gap}\\
\midrule
Li et al. 2017 / 2022~\cite{li2017mdsn,li2022diffusion} & Introduces diffusion auction model, network VCG, IDM, and critical diffusion mechanisms. & Single item, social graph, no diffusion cost, standard quasilinear utilities. & IDM is robust and implementable, but does not target prior-dependent revenue optimality and can sacrifice efficiency.\\
\addlinespace
Zhao et al. 2018~\cite{zhao2018multi} & Multiple units via social networks. & Homogeneous multiple items. & Multi-item version is too large for this project, but useful for future extension.\\
\addlinespace
Zhang et al. 2020 redistribution~\cite{zhang2020redistribution} & Redistribution mechanism on networks; shows direct redistribution is hard in network settings. & Resource allocation with network invitations. & Connects diffusion auctions to budget redistribution; can inspire reward smoothing.\\
\addlinespace
Zhang et al. 2020 fair rewards~\cite{zhang2020fair} & Rewards weak as well as strong critical ancestors. & Single item, no diffusion costs. & Fairness-revenue frontier remains a good experimental and theoretical target.\\
\addlinespace
Bhattacharyya et al. 2024~\cite{bhattacharyya2024optimal} & Myerson-like characterization; revenue-optimal referral auction for i.i.d. bidders. & Referral auction class, i.i.d. MHR/regular style assumptions. & Non-i.i.d. case is explicitly challenging; promising main direction.\\
\addlinespace
Li et al. 2024 transaction costs~\cite{li2024cost} & Adds transaction costs to diffusion auction design. & Cost model for information propagation. & Private heterogeneous edge costs plus non-i.i.d. values are still a hard mixed setting.\\
\addlinespace
Jia et al. 2023~\cite{jia2023dp} & Differentially private diffusion mechanisms. & Single item, privacy constraints. & Privacy adds noise; interaction with optimal referral rewards is open-ish and experimentally testable.\\
\addlinespace
Chen et al. 2023~\cite{chen2023sybil} & Sybil-proof diffusion auctions. & Fake identities allowed. & Strong impossibilities suggest restricted graph classes or approximate goals.\\
\addlinespace
Gu et al. 2023~\cite{gu2023redistribution} & Black-box redistribution framework for truthful diffusion auctions. & Single item, truthful base auction. & Could be combined with non-i.i.d. revenue tuning.\\
\bottomrule
\end{longtable}

\section{Most Promising Research Direction}

\subsection{Main direction: non-i.i.d. robust referral auctions}

Bhattacharyya et al.~develop a Myerson-like viewpoint for referral auctions and obtain revenue-optimal structure for i.i.d. bidders. Their own discussion indicates that moving away from i.i.d. values makes revenue maximization much less structured; in particular, known reductions to threshold auctions may not capture all optimal diffusion auctions.

\begin{researchq}[Non-i.i.d. robustness]
For tree networks with independent but non-identical value distributions, can we characterize a small class of diffusion incentive-compatible mechanisms that strictly dominate IDM in expected revenue while preserving IR and WBB?
\end{researchq}

A realistic course-project version should restrict the graph family. The most tractable choices are:
\begin{enumerate}[itemsep=0.25em]
    \item a path graph,
    \item a rooted tree of height two,
    \item a balanced binary tree with depth-biased value distributions,
    \item a tree plus a small number of cross edges.
\end{enumerate}

The central reason this is promising is that a path or tree preserves enough structure to write threshold conditions, while non-i.i.d. valuations already move beyond the solved i.i.d. optimal-referral theorem.

\subsection{Candidate theorem target}

Consider a rooted path $s-1-2-\cdots-k$ where bidder $i$ has distribution $F_i$. Let a mechanism be \emph{path-threshold referral} if allocation is determined by node-specific monotone thresholds and rewards are paid only to ancestors of the winner.

\begin{conjecture}[Path threshold sufficiency under regularity]
For independent regular distributions on a path, among path-threshold referral mechanisms satisfying IR, DIC, and WBB, expected revenue maximization reduces to optimizing monotone virtual thresholds along the path.
\end{conjecture}

This is not claiming global optimality over all diffusion auctions. It is a deliberately narrow statement. A successful project can either prove the conjecture for $k=2$ or $k=3$, or find a counterexample. Either result is meaningful.

\subsection{Why a counterexample would still be valuable}

In mechanism design projects, a clean impossibility or counterexample is often better than an unproven broad claim. A useful counterexample would show that a natural Myerson-style threshold rule fails under one of:
\begin{enumerate}[itemsep=0.25em]
    \item non-identical distributions,
    \item non-regular distributions,
    \item topology asymmetry,
    \item reward fairness constraints,
    \item private diffusion costs.
\end{enumerate}

\section{Other Theory Directions}

\subsection{Fairness-revenue frontier}

IDM rewards only strong critical ancestors. Fair Diffusion Mechanism argues that weak critical ancestors also contribute to connecting the winner. This suggests a constrained optimization problem:
\[
    \max_M \; \mathbb{E}[\Rev(M)] \quad\text{s.t.}\quad M\text{ is IR, DIC, WBB, and }\Phi(M)\le \tau,
\]
where $\Phi(M)$ is a reward inequality measure such as Gini coefficient or max-min reward gap among contributing nodes.

\begin{researchq}[Reward fairness]
Can one design a family of mechanisms interpolating between IDM-style revenue and FDM-style fairness, while preserving DIC and WBB on trees?
\end{researchq}

A possible small theorem is to show that any reward to a weak critical ancestor must be funded by reducing either seller revenue or strong-ancestor reward, unless the allocation rule changes.

\subsection{Private diffusion costs}

If inviting neighbours has cost, then ``invite all neighbours'' is no longer free. The literature has begun studying transaction costs, but a hard extension is where each edge cost is private and heterogeneous.

\begin{researchq}[Private edge costs]
Suppose buyer $i$ incurs private cost $c_{ij}$ for inviting neighbour $j$. Can any deterministic single-item mechanism be IR, DIC in values and invitations, WBB, and guarantee revenue at least local Vickrey for all cost profiles?
\end{researchq}

A likely outcome is an impossibility theorem. Even a two-hop path can be enough: to make node $1$ invite node $2$, the mechanism may need to compensate node $1$ for both competitive loss and invitation cost. If the seller must remain WBB for all values, compensation may be impossible when the downstream value is low.

\subsection{Topology-sensitive efficiency loss}

IDM's welfare loss is caused by priority along the critical sequence. In dense graphs, critical ancestors are rarer; in tree-like graphs, they are common. This suggests a graph-theoretic bound.

\begin{researchq}[Topology bound]
Can the welfare loss of IDM relative to central Vickrey be bounded using the length of the highest bidder's critical sequence, articulation structure, or dominator-tree depth?
\end{researchq}

A concrete claim to test:
\[
    \frac{\SW(IDM)}{\SW(Central)} \ge f(L),
\]
where $L$ is the length of the diffusion-critical sequence of the global highest bidder. The worst-case bound may be zero, but average-case bounds under random graphs may be informative.

\subsection{Sybil robustness under restricted networks}

Sybil-proof diffusion auctions have strong impossibility phenomena. A smaller project direction is to restrict the graph model: for example, verified institutional networks, bounded-degree trees, or networks where fake identities cannot create arbitrary new edges.

\begin{researchq}[Restricted Sybil model]
On rooted trees where a Sybil attack can only subdivide existing edges, can we design a mechanism that preserves a constant fraction of IDM revenue while being Sybil-proof?
\end{researchq}

This is likely too ambitious as the main project, but it is a strong backup discussion point.

\section{Proposed Project Plan}

\subsection{Minimum publishable-style goal for the course}

The project should aim for one of the following outcomes:
\begin{enumerate}[label=\textbf{G\arabic*.},itemsep=0.35em]
    \item \textbf{A theorem in a restricted model.} Example: characterize optimal path-threshold referral auctions for a two-level tree with independent non-identical regular distributions.
    \item \textbf{A counterexample.} Example: show that a natural extension of i.i.d. optimal referral auctions fails monotonicity or WBB under non-i.i.d. distributions.
    \item \textbf{A robust experimental negative result.} Example: show that revenue gains of tuned referral mechanisms are highly topology-dependent and vanish on dense graphs or under costly diffusion.
\end{enumerate}

\subsection{Experiment design}

The simulator should compare mechanisms under the following axes:
\begin{itemize}[itemsep=0.25em]
    \item \textbf{Topologies:} line, star, balanced tree, Erdos-Renyi, Barabasi-Albert.
    \item \textbf{Valuations:} i.i.d. uniform, lognormal, depth-biased non-i.i.d., community-biased non-i.i.d.
    \item \textbf{Diffusion behavior:} full diffusion, no diffusion, probabilistic diffusion, custom strategic withholding.
    \item \textbf{Metrics:} seller revenue, welfare, participation count, diffusion depth, reward inequality, deficit probability.
\end{itemize}

\subsection{Mechanisms implemented in the accompanying simulator}

The code package implements:
\begin{enumerate}[itemsep=0.25em]
    \item central Vickrey auction,
    \item local Vickrey auction,
    \item network VCG,
    \item IDM,
    \item a tunable referral-auction sandbox with node-specific reserves and referral-share rewards.
\end{enumerate}

The fifth mechanism is intentionally labelled a sandbox, not a theorem-level mechanism. It is useful for searching for patterns and counterexamples under non-i.i.d. distributions.

\section{Expected Claims and How to Validate Them}

\subsection{Claim 1: VCG deficit on sparse chains}

On a path $s-1-\cdots-k$ with $v_k=1$ and all other values $0$, network VCG allocates to $k$ and rewards every intermediate critical node. Seller revenue is negative. The simulator includes this as a unit test.

\subsection{Claim 2: IDM protects seller revenue}

On the same path, IDM remains weakly budget balanced. This should be tested across random values and random trees.

\subsection{Claim 3: non-i.i.d. values create room for prior-dependent mechanisms}

If deeper bidders are stochastically stronger, local Vickrey loses welfare, while IDM may still allocate to an ancestor because of critical-sequence priority. A prior-aware referral mechanism with node-specific reserves may improve expected revenue, but it must be checked against DIC/WBB constraints before being presented as a mechanism theorem.

\subsection{Claim 4: dense graphs reduce critical rewards}

In dense graphs, many buyers have multiple paths from the seller; hence fewer nodes are diffusion-critical. We expect network VCG deficits and IDM rewards to be smaller in dense random graphs than in trees or paths.

\section{Suggested Final Report Structure}

For the final submission, use the following structure:
\begin{enumerate}[itemsep=0.25em]
    \item Problem definition and motivation.
    \item Formal model and definitions: IR, DIC, WBB, diffusion-critical sequence.
    \item Literature survey table.
    \item One selected research question.
    \item Theorem/counterexample/proof attempt.
    \item Simulator design and validation tests.
    \item Experiments and plots.
    \item Discussion: what is solved, what remains open.
\end{enumerate}

\section{Implementation Notes}

The code is organized as follows:

\begin{lstlisting}[basicstyle=\ttfamily\small,frame=single]
simulator/
  refauc/
    instance.py          # model, diffusion, critical sets
    generators.py        # graph and valuation generators
    mechanisms/
      benchmarks.py      # central and local Vickrey
      network_vcg.py
      idm.py
      param_referral.py  # experimental sandbox
    experiments.py
  run_experiments.py
  plot_results.py
  tests/test_mechanisms.py
\end{lstlisting}

The most important implementation detail is the computation of $d_i$. The simulator removes node $i$ from the directed invitation graph and checks which buyers become unreachable from the seller. This directly implements the definition of diffusion-critical descendants.

\section{Risks and Mitigations}

\begin{tabularx}{\textwidth}{p{0.25\textwidth}X}
\toprule
\textbf{Risk} & \textbf{Mitigation}\\
\midrule
Too much literature, too little result & Fix one narrow target: non-i.i.d. path or two-level tree.\\
Mechanism not actually truthful & Clearly separate theorem-level mechanisms from experimental sandboxes.\\
Experiments look like random plotting & Every plot should test a specific claim: deficit, welfare loss, revenue gain, or topology effect.\\
Open problem too hard & Aim for a counterexample or restricted theorem instead of broad optimality.\\
\bottomrule
\end{tabularx}

\section{Conclusion}

The best version of this project is theory-first but implementation-supported. The literature already contains strong mechanisms for single-item diffusion auctions, fairer rewards, redistribution, privacy, transaction costs, and i.i.d. optimal referral auctions. Therefore, the project should focus on a boundary of known theory: non-i.i.d. values and topology robustness. A small theorem, a clean counterexample, or a carefully validated experimental insight would align well with a research-style mechanism design project.

\begin{thebibliography}{99}

\bibitem{li2017mdsn}
Bin Li, Dong Hao, Dengji Zhao, and Tao Zhou.
\newblock Mechanism Design in Social Networks.
\newblock \emph{AAAI}, 2017.

\bibitem{li2022diffusion}
Bin Li, Dong Hao, Hui Gao, and Dengji Zhao.
\newblock Diffusion auction design.
\newblock \emph{Artificial Intelligence}, 303:103631, 2022.

\bibitem{zhao2018multi}
Dengji Zhao, Bin Li, Junping Xu, Dong Hao, and Nicholas R. Jennings.
\newblock Selling Multiple Items via Social Networks.
\newblock \emph{AAMAS}, 2018.

\bibitem{zhang2020redistribution}
Wen Zhang, Dengji Zhao, and Hanyu Chen.
\newblock Redistribution Mechanism on Networks.
\newblock \emph{AAMAS}, 2020.

\bibitem{zhang2020fair}
Wen Zhang, Dengji Zhao, and Yao Zhang.
\newblock Incentivize Diffusion with Fair Rewards.
\newblock \emph{ECAI}, 2020.

\bibitem{guo2021survey}
Yuhang Guo and Dong Hao.
\newblock Emerging Methods of Auction Design in Social Networks.
\newblock \emph{IJCAI Survey Track}, 2021.

\bibitem{bhattacharyya2024optimal}
Rangeet Bhattacharyya, Parvik Dave, Palash Dey, and Swaprava Nath.
\newblock Optimal Referral Auction Design.
\newblock \emph{AAMAS}, 2024.

\bibitem{li2024cost}
Bin Li, Dong Hao, and Dengji Zhao.
\newblock Diffusion auction design with transaction costs.
\newblock \emph{Autonomous Agents and Multi-Agent Systems}, 38(2), 2024.

\bibitem{jia2023dp}
Fengjuan Jia, Mengxiao Zhang, Jiamou Liu, and Bakh Khoussainov.
\newblock Incentivising Diffusion while Preserving Differential Privacy.
\newblock \emph{UAI}, 2023.

\bibitem{chen2023sybil}
Hongyin Chen, Xiaotie Deng, Ying Wang, Yue Wu, and Dengji Zhao.
\newblock Sybil-Proof Diffusion Auction in Social Networks.
\newblock \emph{AAMAS}, 2023.

\bibitem{gu2023redistribution}
Sizhe Gu, Yao Zhang, Yida Zhao, and Dengji Zhao.
\newblock A Redistribution Framework for Diffusion Auctions.
\newblock \emph{AAMAS}, 2023.

\bibitem{myerson1981}
Roger B. Myerson.
\newblock Optimal Auction Design.
\newblock \emph{Mathematics of Operations Research}, 6(1):58--73, 1981.

\end{thebibliography}

\end{document}
